\color{unidarkgreen}{\subsection{Максимальная токовая защита стороны ВН}}
\color{black}

\begin{enumerate}

\item
Максимальная токовая защита (МТЗ) стороны ВН предназначена для резервирования действия основных быстродействующих защит. Защита выполнена двухступенчатой. Измерительные цепи защиты включены на токи фаз А, В, С со стороны ВН трансформатора.    

\item
Каждая из ступеней имеет регулируемые значения по току и времени срабатывания. Значения срабатывания по току определяются параметрами \textbf{«Iср МТЗ-1 ВН»} и \textbf{«Iср МТЗ-2 ВН»} для первой и
второй ступеней соответственно. Задержки на срабатывание защиты задаются параметрами \textbf{«Тср МТЗ-1 ВН»} и \textbf{«Тср МТЗ-2 ВН»}.

\item
Каждая из ступеней защиты может выполнена с пуском по напряжению. Ввод контроля пуска по напряжению в логику работы ступеней осуществляется программными накладками \textbf{«Режим КПОН МТЗ-1 ВН»} и \textbf{«Режим КПОН МТЗ-2 ВН»} для первой и второй ступеней соответственно. 

\item
Каждая из ступеней защиты может выполнена с блокировкой по второй гармонике (отношению уровня второй гармоники к первой). Данная блокировка реализована на базе модуля «ОВ БНТ ВН». Ввод контроля броска намагничивающего тока программными накладками \textbf{«Режим БНТ МТЗ-1 ВН»}, \textbf{«Режим БНТ МТЗ-2 ВН»}.

\item
Для обеспечения несрабатывания МТЗ ВН при броске тока намагничивания используется перекрестная блокировка по 2 гармонике в фазных токах ВН. Вводится программной накладкой \textbf{«БНТ МТЗ ВН»}. Блокировка основана на контроле отношения действующего значения второй гармонике к первой токов ВН. Данная блокировка вводиться в работу для каждой ступени независимо с помощью уставок \textbf{«Режим БНТ МТЗ-1 ВН»} и \textbf{«Режим БНТ МТЗ-2 ВН»}. Имеется возможность регулировать порог срабатывания блокировки, с помощью уставки \textbf{«I2г/I1г МТЗ ВН»}. Блокировка МТЗ ВН по 2 гармонике выводится из действия при величине соответствующего фазного тока ниже 0,04Iном. Также блокировка МТЗ ВН по 2 гармонике выводится из действия при достижении величины действующего значения тока выше уставки \textbf{«Iср БНТ МТЗ ВН»}.

\item
Алгоритм работы МТЗ ВН приведен на рисунке \ref{pic:MTZ_VN}. Параметры для настройки МТЗ ВН приведены в таблице \ref{app_set:mtzvn}.

\medskip
\begin{figure}[H]
\centering
\includegraphics[width=0.8\textwidth,height=0.8\textheight,keepaspectratio]{\fbpath/06.01 ДЗТ/МТЗ ВН/img/MTZ_VN.pdf}
\captionof{figure}{Функциональная схема логики МТЗ ВН}\label{pic:MTZ_VN}
\end{figure}

\end{enumerate}


